\chapter{Complex Paths and Analytic Continuation}
\label{ch:complex-paths}

Complex analysis can begin with polynomial identities, finite paths, and
the series charts just constructed.  The contour, its subdivision, and
the local bounds are supplied data.  We prove the integral identities
actually used, rather than first postulating an arbitrary holomorphic
function on an arbitrary open set.

\section{Integrating on a segment}
For rational complex endpoints $p,q$, define
\[
 \int_{[p,q]}F(z)\,dz
 =(q-p)\int_0^1 F(p+t(q-p))\,dt.
\]
The right side consists of two real integral computations.  A polygonal
path is a finite list of such segments, and its integral is their sum.
Its length has a rational upper bound obtained from coordinate differences.
If two integrands differ by at most $\eta$ on the path of length at most
$L$, their integrals differ by at most $L\eta$.

In a single series disk, the primitive
\[
 P(c+w)=\sum_{n\ge0}\frac{a_n}{n+1}w^{n+1}
\]
has a geometric tail on every smaller disk.  Its finite polynomial
identities and their tails prove
\[
 \int_{[p,q]}F=P(q)-P(p)
\]
whenever the segment stays in such a disk.  This also holds for a
specified continuously differentiable parametrized arc in the disk: the
finite polynomial chain rule gives the result, and its error is bounded
by a supplied arc-length bound.  Circle arcs from the earlier geometric
construction are one such family.

\section{A finite form of the contour theorem}
\label{thm:finite-contour}
Suppose a rational triangulated polygon is supplied, and each closed
triangle lies in a smaller certified series disk of $F$.  The descriptions
are proved to agree on their overlaps.  The integral around each triangle
is zero, by its local primitive.  Add these finitely many identities:
interior edges cancel with their opposite orientations.  Therefore
\[
 \int_{\partial\Omega}F(z)\,dz=0.
\]
A polygon with holes is treated by the same cancellation, giving the
outer boundary integral as the sum of the positively oriented inner
boundary integrals.  A finite chain of these subdivisions proves equality
of two path integrals.  No existence of a finite subcover is assumed:
the finite subdivision and the certified disk margins are exhibited.

This is also a computational form of contour deformation.  It is
important to retain the overlap comparisons.  Merely evaluating the
same printed formula on two sides of a branch cut is not a proof that
one has continued the same function.

\section{Coefficient extraction on a circle}
For $0<r<R$, parametrize the circle by
\[
 z(t)=c+rE(2\pi it),\qquad 0\le t\le1.
\]
This is computed at rational $t$ by trigonometry, not by assuming circle
points are rational.  Its derivative and length bound have already been
proved.  Finite trigonometric orthogonality, or the trigonometric FTC,
gives for every integer $j$
\[
 \int_0^1 E(2\pi ijt)\,dt=\begin{cases}1,&j=0,\\0,&j\ne0.\end{cases}
\]
Integrate a finite polynomial and use its uniform series tail.  The result is
\[
 \boxed{a_k=\frac1{2\pi i}\int_{|z-c|=r}
                  \frac{F(z)}{(z-c)^{k+1}}\,dz.}
\]
If the value boxes on this circle certify $|F|\le B$, this formula gives
$|a_k|\le B r^{-k}$.  Thus an independent boundary estimate can produce
better coefficient bounds, rather than just use the original majorant.

For a computed $w$ with $|w-c|<r$, expand
$(z-w)^{-1}=\sum_{j\ge0}(w-c)^j/(z-c)^{j+1}$ on the circle.  Its tail
is bounded by a geometric series with ratio $|w-c|/r$.  Finite integration
and the coefficient formula give the local Cauchy identity
\[
 \boxed{F(w)=\frac1{2\pi i}\int_{|z-c|=r}
                         \frac{F(z)}{z-w}\,dz.}
\]
Both sides are now specified computations with independent truncation and
quadrature errors.  We have proved the identity for the supplied chart,
not invoked it to manufacture a chart from unverified pointwise data.

\section{Arctangent normalizes the Cauchy pole}
\label{sec:cauchy-arctan}
There is a direct rectangle route to the same normalization. Split the
counterclockwise boundary of $[-1,1]^2$ into eight affine half-edges.
On the right side, $z=1+iu$ gives
\[
 \frac{dz}{z}=\frac{u+i}{1+u^2}\,du.
\]
Pair opposite parameter signs and rotate through the four sides. The real
parts cancel exactly, leaving $8i/(1+u^2)$ on $[0,1]$. Our rational boxes
enclose every tagged sum on the common partition and have imaginary width
at most $16/2^n$. Thus the checked square computation satisfies
\[
 \oint_{\partial[-1,1]^2}\frac{dz}{z}=8iA(1)=2\pi i,
\]
with the original geometric $\pi$. This is
\texttt{PDE.CauchyContour.raw\_equiv\_twoPiI}; translation and rational
nonzero dilation preserve its pole pullback.

PNT+ formalizes the corresponding rectangle argument in Mathlib:
arctangent and logarithm side primitives normalize the pole, while
Mathlib's rectangle Cauchy theorem removes the regular part. Its checked
finite residue theorem concerns simple poles with no boundary poles.
Applying that theorem to
\[
 \frac{F(z)}{z-p}=
 \frac{F(z)-F(p)}{z-p}+\frac{F(p)}{z-p}
\]
gives the Cauchy value formula once the divided difference is extended
holomorphically across $p$. PNT+'s explicitly named derivative formula
uses a separate Mathlib circle-integral theorem. These are external
formal references, not imports into our foundation.

The native square normalization is checked. General native contour
transport, regular-part removal and the Cauchy integral formula still
require their explicit subdivision, local chart and convergence bridges.
The adjacent general contour descriptions give the mathematical program,
not additional checked declarations supplied by this normalization.
See the source ledger \texttt{docs/PNT\_CAUCHY\_ARCTAN.md}.

\section{Known poles and finite residues}
Suppose near a specified point $s$ the computation has a proved description
\[
 F(z)=\sum_{j=1}^m b_{-j}(z-s)^{-j}+H(z),
\]
where $H$ has a series chart.  On a sufficiently small certified circle,
finite coefficient extraction gives
\[
 \int F(z)\,dz=2\pi i b_{-1}.
\]
The coefficient $b_{-1}$ is called the residue.  The finite contour
cancellation above transports this result to larger supplied polygonal
contours when the region between them is covered by certified charts.
For circular boundary components, approximate the arcs by the geometric
polygonal construction inside the annular chart cover.  Uniform continuity
and path-length bounds control the integral errors; radial joining edges
cancel.  This supplies the same transport for those boundaries.

As a direct example, on $|z|=2$,
\[
 \frac1{z(z-1)}=\sum_{n\ge0}z^{-n-2},
\]
so its integral is zero.  Around $0$ and $1$ separately the residues are
$-1$ and $1$.  The two ways of computing the contour integral agree.
The poles were supplied and their expansions proved; no unsupported
search for every pole or zero has taken place.

\section{Logarithm along a specified path}
Start near $1$ with
\[
 L(1+w)=\sum_{n\ge1}\frac{(-1)^{n+1}}n w^n,\qquad |w|<1.
\]
Its derivative is the geometric series for $1/(1+w)$.  The series product
and chain rules show $E(L(1+w))=1+w$: its quotient by $1+w$ has all
nonconstant derivative coefficients zero and constant value $1$.
On the positive real interval it agrees with the previously constructed
logarithm by its primitive identity and the value at $1$.

Suppose a value $L(c)$ has been continued to a nonzero computed point $c$.
Choose a certified $d<|c|$ and set
\[
 L(c+w)=L(c)+\sum_{n\ge1}\frac{(-1)^{n+1}}n(w/c)^n,
 \qquad |w|\le d.
\]
This is a fresh chart, whose radius comes from separation from zero.
The value at the next center is computed from the old chart, and the two
charts have derivative $1/z$ on their overlap.  Their local primitive
identities and matching value prove agreement on the overlap, which is
convex.  A finite chain of such steps computes continuation along a
supplied path avoiding zero.  A rational positive distance from zero on
each path piece gives a finite step schedule.

The end value can depend on the path.  On the unit circle,
\[
 \int\frac{dz}{z}=\int_0^1 2\pi i\,dt=2\pi i.
\]
Thus continuation from $1$ once counterclockwise changes the logarithm
value by $2\pi i$.  Starting with value $0$, the computation $E(L/2)$
returns to $-1$, not $1$.  This constructs the sign change of the square
root after a circuit.  A path and a starting branch are part of the input;
there is no single globally chosen square-root computation on the
punctured plane.

\section{When does a chain of charts count as continuation?}
A continuation certificate consists of finitely many local charts, steps
lying strictly within their certified disks, and proved comparisons on
each successive overlap.  Comparisons can be coefficient identities,
local primitive identities, or the uniqueness estimates for an explicitly
constructed differential or integral equation.  Equality of all Taylor
coefficients at a common center implies equality on their common smaller
disk by the finite tail estimates.  Matching only a finite jet is not
enough unless an equation proves that it determines the later coefficients.

A path change is justified by a finite subdivision whose cells stay in
such common charts.  It does not follow from drawing apparently similar
paths.  This requirement exposes the separation and branch information
that an appeal to unrestricted analytic continuation would hide.

The next chapters use two different ways to replenish local bounds:
parameter integrals supply new majorants in right half-planes, while
linear equations supply new evolution bounds away from their coefficient
singularities.  Both can be connected to these local contour computations.
